% Generated draft: requires mathematical and source review.
\documentclass[11pt]{article}
\usepackage[margin=25mm]{geometry}
\usepackage{fontspec}
\setmainfont{Latin Modern Roman}
\usepackage{amsmath,amssymb,mathtools,unicode-math}
\setmathfont{Latin Modern Math}
\usepackage{enumitem,needspace,xurl,hyperref}
\setlength{\parindent}{0pt}
\setlength{\parskip}{5pt}
\setlength{\emergencystretch}{3em}
\newcommand{\N}{\mathbb{N}}
\newcommand{\Z}{\mathbb{Z}}
\newcommand{\Q}{\mathbb{Q}}
\newcommand{\R}{\mathbb{R}}
\newcommand{\C}{\mathbb{C}}
\newcommand{\F}{\mathbb{F}}
\newcommand{\A}{\mathbb{A}}
\newcommand{\PP}{\mathbb{P}}
\newcommand{\E}{\mathbb{E}}
\newcommand{\eps}{\varepsilon}
\newcommand{\dd}{\,\mathrm{d}}
\begin{document}
\section*{UnsolvedMath ID: 30005931}
\textbf{Source problem code:} catalogue entry
\subsection{Definitions and mathematical statement}
For terminology only, let \(G=(V,E)\) be a graph equipped with a percolation configuration \(\omega\), with openness of paths interpreted according to the chosen percolation model. For vertices \(x,y\) admitted by that configuration, the chemical distance is
\[
d_\omega(x,y)=\inf\bigl\{n\in\Z_{\geq 0}:\text{there is an open path of }n\text{ edges from }x\text{ to }y\bigr\},
\]
with \(\inf\varnothing=+\infty\) and the length-zero path allowed at an admitted vertex. This definition explains the terminology; the supplied record does not specify the intended graph, percolation model, critical parameter, or conditioning on connectivity.

\textbf{Source research request.} Obtain stronger results for chemical distance in critical percolation and its relation to metrics associated with nonsimple conformal loop ensembles (CLE).

\textbf{Unresolved formulation.} The supplied record does not specify the existing result to be strengthened, a quantitative bound or convergence assertion, the CLE parameter and metric space, or the intended relation between the percolation and CLE metrics. Consequently, it does not determine a precise quantified conjecture. Its dated literature assessment mentions an exact chemical-distance exponent and a stronger identification as remaining targets, but supplies neither a definition of that exponent in a specified model nor the identification to be proved. These omissions must be resolved from reliable source material before a precise theorem-shaped statement can be given. No particular lattice, domain, boundary condition, distance normalization, exponent value, or mode of convergence is imposed in this draft.
\subsection{Short English statement}
The record asks for stronger results about shortest open-path distances in critical percolation and their relation to nonsimple CLE metrics. It leaves the model, comparison theorem, and desired strengthening unspecified, so this is a research request requiring source clarification rather than a fully quantified conjecture. The supplied designation \emph{partially solved} is dated curation, not an independently verified mathematical conclusion.
\subsection{Sources}
\begin{itemize}
\item Supplied local UnsolvedMath registry record, permanent ID \(30005931\), title \emph{Chemical Distance in Critical Percolation}, source problem code \texttt{catalogue entry}; collection: \emph{Oberwolfach Reports Open Problems}; proposed year: 2024.
\item \emph{Statistical Physics and Random Surfaces} (2024), Oberwolfach Reports, pp.\ 1387--1444, DOI \texttt{10.4171/owr/2024/25}. Supplied source report and literature reference labelled ``OWR 2024/25'': \url{https://doi.org/10.4171/owr/2024/25}. The supplied background describes these reports as documenting workshops and their open-problem discussions. No precise passage or individual problem page is supplied.
\item Attribution qualification: the original extraction contains the incomplete phrase ``This problem was asked by O.'' followed by ``The chemical distance metric for non-simple CLE Yizheng Yuan (joint work with Valeria Ambrosio and Jason Miller)''. This fragment associates Yizheng Yuan, Valeria Ambrosio, and Jason Miller with the surrounding material but does not reliably identify the proposer or establish where the problem text ends and the adjacent contribution begins. The registry's proposer field is empty.
\item Supplied literature reference labelled ``Intrinsic metric scaling limit'': \url{https://arxiv.org/abs/2604.14122}. No authors or detailed theorem statement are supplied for this reference.
\item Supplied literature reference labelled ``Damron--Hanson--Sosoe'': \url{https://arxiv.org/abs/1708.03643}.
\item Literature assessment checked at \texttt{2026-08-22}: the supplied curation labels the problem \emph{partially solved} and states that bounds and a 2026 intrinsic-metric scaling-limit result substantially advance critical chemical distance, while the exact chemical-distance exponent and a stronger identification remain open. This is supplied dated curation derived from the OWR material, not independently verified here; its time-sensitive status and cited claims require independent verification.
\item The registry labels its statement \texttt{corrected\_verified} and says that only formula delimiters, notation, grammar, and typography were repaired while preserving mathematical content and scope. This supplied label does not resolve the missing hypotheses or the damaged attribution fragment.
\end{itemize}
\end{document}
